% ============================================================================
% lim_section.tex -- R4 deliverable: the lim subsection, paper form.
% Placement: Section 5 (Variants), immediately after "Restarting from Zero"
% (i.e. right before \subsection{Both Directions at Once}).  Every numeric
% claim is machine-verified; the scripts and domains are named in the
% (Verified: ...) notes and in research/scratch/lim/REPORT.md (R2-R4):
%   verify_r2.py (round2.log)   -- the compilation, end to end
%   verify_r3.py (round3.log)    -- occurrence census, normalizations, driver
%   verify_r4.py  (round4.log)  -- every number this section states
% Compiles standalone: 0 errors (see the wrapper in the sanity harness).
% ============================================================================

\subsection{Converging to the Fixed Point}
\label{ssec:lim}

Restart iterates a \emph{rule}; we now iterate an \emph{expression}. The
variant of this subsection adds one node to the grammar: apply the
unary expression to its own output until nothing changes. Where
Definition~\ref{def:markov} normalizes one rewrite rule
$B \to A$ under the leftmost strategy, the new node normalizes a whole
pipeline under the strategy the calculus already has -- the sweep -- and
the two normalizers of the same rule will turn out to be genuinely
different functions (Remark~\ref{rem:limstrategies} below). The reward for the
difference is a characterization the restart calculus does not have:
with this one node the flat calculus computes \emph{exactly} the partial
computable functions, and the entire distance between the bounded
calculus of Theorem~\ref{thm:fp} and universality turns out to be one
iteration node (Corollary~\ref{cor:limclass}).

\begin{definition}
    (\emph{Fixed-Point Iteration})
    \label{def:lim}
    Extend the grammar of Definition~\ref{def:exp} by the rule
    \begin{mathpar}
        \inferrule{E \in \mathrm{Exp}_1, \; E_0 \in \mathrm{Exp}_n}{\mathrm{lim}(E)[E_0] \in \mathrm{Exp}_n}
    \end{mathpar}
    where $X_1$ is \emph{bound} in $E$ by the node; the free variables of
    $\mathrm{lim}(E)[E_0]$ are those of $E_0$. The \emph{orbit} of
    $E \in \mathrm{Exp}_1$ on $w$ is $s_0 = w$ and
    $s_{k+1} = \llbracket E \rrbracket(s_k)$, and
    \begin{align}
        \llbracket \mathrm{lim}(E)[E_0] \rrbracket(\vec S) \;=\; s_K
        \quad \text{for the least $K$ with } s_{K+1} = s_K\,,
    \end{align}
    with $s_0 = \llbracket E_0 \rrbracket(\vec S)$: the first state the
    iterand does not change. If some $s_{k+1}$ is undefined, or no such
    $K$ exists, the node is undefined (in the second case it
    \emph{diverges}). Substitution acts on the start expression,
    \begin{align}
        \mathrm{lim}(E)[E_0]\big[F_1/X_1, \dots, F_n/X_n\big]
        \;=\; \mathrm{lim}(E)\big[\,E_0[F_1/X_1, \dots, F_n/X_n]\,\big],
    \end{align}
    capture-free, since only $X_1$ is bound and only in $E$ -- so
    Lemma~\ref{lem:beta} survives verbatim for the extended grammar.
\end{definition}

An $L{+}\mathrm{lim}$ expression is \emph{flat} if its iterands contain no
$\mathrm{lim}$ nodes; a single node $\mathrm{lim}([\,A/B\,])$ iterates one
constant pass, the \emph{sweep} $w \mapsto [\,A/B\,]w$. The preceding
remark's sort has its sweep twin: $\mathrm{lim}([\,ba/ab\,])$ also maps
every input to $b^{\#b}a^{\#a}$ (verified on all binary inputs of length
$\leq 12$ and $100$ random ones of length $13$--$16$).

\begin{lemma}
    (\emph{The Fixed Points of a Pass})
    \label{lem:fpsweep}
    For $B \neq \epsilon$ and $A \neq B$:
    \begin{align}
        [\,A/B\,]w = w \quad\Longleftrightarrow\quad B \not\subset w .
    \end{align}
\end{lemma}

\begin{proof}
    If $B$ does not occur, the scan never fires. If it does, the sweep
    performs $m \geq 1$ non-overlapping replacements, changing the length
    by $m(|A| - |B|)$ -- so $w$ is changed whenever $|A| \neq |B|$ -- and
    when $|A| = |B|$ (with $A \neq B$) the first replacement spells $A$
    where $w$ spells $B$ at the same position, with the prefix before it
    untouched.
\end{proof}

So for $A \neq B$ the sweep's fixed points are exactly the $B$-free
strings -- the same stop set as the restart process of
Definition~\ref{def:markov}, which stops at the first $B$-free
trajectory term. Lemma~\ref{lem:fpsweep} is also a complete decision
procedure for the \emph{occurrence} of a constant, and at one stroke a
subroutine of the compilation below: a pass with a one-character
replacement tests containment,
\begin{align}
    \mathtt{contains}_B(E) \;\triangleq\; \mathtt{eq}\big([\,c/B\,]E,\; E\big),
    \qquad c \text{ any character with } \{c\} \neq B,
    \label{eq:contains}
\end{align}
because $[\,c/B\,]E = E$ iff $B \not\subset E$.

\begin{proposition}
    (\emph{One-Rule Totality})
    \label{prop:limtotal}
    Let $B \neq \epsilon$.
    \begin{enumerate}[(i)]
        \item $|A| \leq |B|$ and $A \neq B$: $\mathrm{lim}([\,A/B\,])$ is total.
        \item $B \subset A$: $\mathrm{lim}([\,A/B\,])$ diverges on every $B$-containing input.
        \item $A = B$: $\mathrm{lim}([\,A/B\,])$ is the identity.
    \end{enumerate}
\end{proposition}

\begin{proof}
    The proofs of Theorem~\ref{thm:termination}(i) and (ii) bound a
    replacement at \emph{any} position -- length in (i), the potential
    $V$ in (ii) -- so each sweep with at least one replacement strictly
    decreases a measure that takes finitely many values, and the orbit
    of sweeps reaches a $B$-free string, a fixed point by
    Lemma~\ref{lem:fpsweep}: (i). For (ii), every replaced site
    re-contains $B$, so after every sweep the string still contains $B$
    and is not a fixed point; for (iii) the pass is the identity, so the
    first sweep already is a fixed point. Note the one divergence from
    the restart column: at $A = B$ the restart process still hunts a
    redex and diverges (Theorem~\ref{thm:termination}(iii)), while the
    sweep stops -- the two stop rules agree exactly when $A \neq B$.
\end{proof}

\begin{remark}
    (\emph{Two Normalizers, Neither Canonical})
    \label{rem:limstrategies}
    $\mathrm{lim}([\,A/B\,])$ and $[\,A/B\,]^{\mathrm{m}}$ are two
    deterministic normalizers of the one-rule system $B \to A$, and
    one-rule systems are not confluent: they disagree. On the census
    domain of Theorem~\ref{thm:termination}(v) -- the $930$ binary rules
    with $|A|, |B| \leq 4$ on all inputs of length $\leq 9$, restart at
    a $3 \cdot 10^5$ step cap -- the two disagree on \emph{value} on
    $3{,}306$ pairs and on \emph{termination} on $290$ pairs, the latter
    exactly the two rules $[\,abba/bab\,]$ and $[\,baab/aba\,]$, all in
    the direction restart-diverges, sweep-converges, and \emph{none} in
    the reverse direction; re-verified on all inputs of length $\leq 6$
    at a $2{,}000$ step cap ($126$ value, $20$ termination disagreements,
    $0$ reverse). Witnesses, at a $10^5$ step restart cap:
    $[\,ab/bb\,]$ on $bbbb$ gives $\mathtt{abab}$ (one sweep replaces
    the two disjoint occurrences) against restart's $\mathtt{aaab}$
    (three sequential leftmost splices, each shifting the next site);
    $[\,ab/abba\,]$ on $abbababba$ gives $abb$ in two sweeps against
    restart's $abbba$; and $[\,baab/aba\,]$ on $aaaba$ goes
    $aaaba \to aabaab \to abaabab \to baabbaabb$ -- three sweeps, done --
    where restart diverges, each leftmost splice re-creating an $aba$ at
    the junction, the census mechanism of
    Theorem~\ref{thm:termination}(v). The sweep is thus
    \emph{sometimes} the better-defined of the two strategies, and
    sometimes not definable from it at all: neither normalizer computes
    the other's values on a non-confluent rule.
\end{remark}

The rest of this subsection shows that this modest operator -- iterate a
flat expression to its fixed point -- is exactly the missing piece
between the polynomial calculus of Theorem~\ref{thm:fp} and universality.

\begin{theorem}
    (\emph{Compiling Two-Counter Machines Flat})
    \label{thm:lim2cm}
    Fix $\Sigma$ with $|\Sigma| \geq 2$ and two characters
    $\sigma \neq \tau \in \Sigma$. For every two-counter machine $M$ with
    instructions $\mathrm{inc}(r, j)$, $\mathrm{decjz}(r, j_0, j_1)$ and
    halt (counters $r \in \{1, 2\}$), there are flat expressions
    $\mathrm{INIT}, \mathrm{step}, \mathrm{OUT} \in \mathrm{Exp}_1$ such
    that
    \begin{align}
        \mathrm{MAIN} \;\triangleq\; \mathrm{OUT}\big[\,F/X_1\,\big],
        \qquad F \;\triangleq\; \mathrm{lim}(\mathrm{step})\big[\,\mathrm{INIT}\,\big],
    \end{align}
    has exactly one $\mathrm{lim}$ node and computes the partial
    function of $M$ on tallies: $\mathrm{MAIN}(\tau^n)$ is the
    counter-2 tally $\tau^{f(n)}$ when $M$ halts on input $n$ with
    counter 2 equal to $f(n)$ and counter 1 empty, and is undefined
    exactly when $M$ does not halt. The iterand $\mathrm{step}$ is an
    if-chain of the Section~\ref{sec:toolkit} constructions over
    constant occurrence tests~(\ref{eq:contains}), all of whose branches
    are constant passes.
\end{theorem}

\begin{proof}[Proof (construction)]
    Three compile-time transformations of $M$ come first, each
    preserving its partial function: (N1) a $\mathrm{decjz}$ whose
    zero-branch jumps to itself would leave the configuration
    untouched -- a non-halting fixed point of the compiled step -- so it
    is replaced by a fresh self-looping increment (both diverge);
    (N2) the output is read from counter 2 (swap the counters if $M$
    leaves its result in counter 1); (N3) jumps to the halt
    instruction are retargeted through a fresh draining instruction
    $\mathrm{decjz}(1, \mathrm{halt}, \mathrm{drain})$, so counter 1 is
    empty at halt.

    \emph{Frames.} Write
    $V(k) = \sigma^2(\tau\sigma)^{k-1}\tau\sigma^3$ for $k \geq 1$:
    $k$ isolated $\tau$'s, head run $\sigma^2$, interior $\sigma$-runs
    of length $1$, tail run $\sigma^3$, length $2k+4$. The compiler
    assigns each \emph{frame constant} a distinct $V(k)$ with $k \geq 2$:
    delimiters $D_0, D_1, D_2, D_3$, instruction codes
    $P_1, \dots, P_s, P_H$ for the instructions of the normalized
    machine, and the wrap marker $M$.

    \emph{Configurations.} For $i \in \{1, \dots, s\}$ let $\Pi(i)$ be
    the program list $P_1 \cdots P_s P_H$ with the $i$-th code wrapped,
    and $\Pi(0) = P_1 \cdots P_s\, M P_H M$ the halted list. The
    configuration with instruction counter $i$ and counter values $x, y$
    is
    \begin{align}
        C(i, x, y) \;=\; D_0\;\Pi(i)\;D_1\;\tau^{x}\;D_2\;\tau^{y}\;D_3 ,
    \end{align}
    an empty tally meaning adjacent delimiters.

    \emph{The step.} On a configuration the current instruction is the
    wrapped one, so dispatch is the occurrence test of
    $M P_i M$~(\ref{eq:contains}): $\mathrm{step}$ is the if-chain
    $\mathtt{if}(\mathtt{contains}_{M P_i M}, \mathrm{body}_i,
    \cdots)$ over $i = 1, \dots, s$, with the identity as the final
    else -- the halted configuration is a fixed point. Each
    $\mathrm{body}_i$ is three constant passes:
    \begin{enumerate}[(a)]
        \item \emph{operate}: for $\mathrm{inc}(r, j)$, the pass
        $[\,D_r\tau/D_r\,]$ (a new tally character at the left edge of
        counter $r$'s run); for $\mathrm{decjz}(r, j_0, j_1)$, the test
        $\mathtt{contains}_{D_r D_{r'}}$ -- adjacent delimiters, i.e.\
        counter $r$ empty (Lemma~\ref{lem:occurrence}(iii)) -- selects
        between the jump $j_0$ and $[\,D_r/D_r\tau\,]$ followed by the
        jump $j_1$ (the decrement deletes the first tally character,
        pinned by Lemma~\ref{lem:occurrence}(iv));
        \item \emph{unwrap}: $[\,P_i/M P_i M\,]$;
        \item \emph{wrap}: $[\,M P_j M/P_j\,]$ for the jump target $j$.
    \end{enumerate}
    Unwrap before wrap: on a self-jump ($j = i$) the wrap pattern
    $P_j$ would otherwise occur inside the block being unwrapped, and
    double-wrap. Every pattern is a nonempty constant, so every branch
    is total on every input, and an unselected branch's passes are
    merely \emph{discarded}: nothing needs the lazy gate of
    Section~\ref{sec:recursion}, because a flat branch has no
    undefined part to guard. (For the doubling machine:
    $\mathrm{step}$ is $298$ nodes, $83$ passes, of which $75$
    constant-pattern; $\mathrm{MAIN}$ is $310$ nodes.)
\end{proof}

\begin{lemma}
    (\emph{Occurrence})
    \label{lem:occurrence}
    Let the frame constants be assigned pairwise distinct family
    indices $k \geq 2$. Then for every instruction counter
    $i \in \{0, \dots, s\}$ and every $x, y \geq 0$, in $C(i, x, y)$:
    \begin{enumerate}[(i)]
        \item every frame constant occurs exactly at its slots ($M$
        twice, at the two wrap positions; every other constant once);
        \item $M P_i M$ occurs iff $i$ is the current instruction
        ($i \geq 1$), and $M P_H M$ iff $i = 0$; in either case exactly
        once;
        \item $D_r D_{r'}$ occurs iff counter $r$ is zero, exactly once;
        \item $D_r\tau$ occurs iff counter $r$ is nonzero, exactly
        once, ending at the \emph{first} tally character;
        \item the output prefix $D_0\,\Pi(0)\,D_1 D_2$ occurs iff
        $i = 0$ and $x = 0$, exactly once, at position $0$.
    \end{enumerate}
\end{lemma}

\begin{proof}
    The proof is a \emph{gap calculus}. For a string $W$ containing at
    least one $\tau$ and no $\tau\tau$, let its \emph{gap sequence} be
    $(h;\, g_1, \dots, g_{m-1};\, t)$: the number of $\sigma$'s before
    the first $\tau$, between consecutive $\tau$'s, and after the last.

    \emph{Step 1 (configurations).} Every frame constant has gap
    sequence $(2;\, 1^{k-1};\, 3)$. Concatenating two constants merges
    the trailing $\sigma^3$ of the first with the leading $\sigma^2$ of
    the second: an interior gap of $5$. A tally $\tau^{n}$, $n \geq 1$,
    between its delimiters contributes $3\,0^{\,n-1}\,2$; an empty tally
    makes them touch, a gap of $5$. So a configuration's gaps are: head
    $2$, then the constants' interior $1$-blocks interleaved with
    junction gaps of $5$, the two tally junctions $3\,0^{x-1}\,2$ or
    $5$, tail $3$.

    \emph{Step 2 (patterns cannot cross tallies).} Every pattern of
    the construction has at least two $\tau$'s, no $\tau\tau$, head
    exactly $2$, tail in $\{0, 3\}$, and interior gaps in $\{1, 3, 5\}$
    only -- never $0$ and never $2$. If a pattern $W$ occurs in
    $C(i,x,y)$, its $\tau$'s map to \emph{consecutive} $\tau$'s of the
    configuration (a skipped $\tau$ would have to appear among $W$'s
    $\sigma$'s), so every interior gap of $W$ equals the corresponding
    gap of the configuration. Gaps of $0$ arise only inside tallies of
    length $\geq 2$, so $W$ covers no two consecutive tally characters;
    a gap of $2$ arises only after the last character of a tally, so
    $W$'s $\tau$-sequence cannot cross a tally's right edge; and if
    $W$'s first $\tau$ were a tally character, its next gap would have
    to be $0$ or $2$ -- both impossible. The only pattern that can meet
    a tally at all is the decrement pattern $D_r\tau$, whose last gap
    $3$ pins its last $\tau$ to the tally's \emph{first} character.
    Every other pattern avoids tallies entirely.

    \emph{Step 3 (pinning).} Let $W$'s leading segment (up to its
    first gap of $3$ or $5$) contain $k$ $\tau$'s, all separated by
    gaps of $1$. Its first $\tau$ must be a constant's \emph{first}
    $\tau$: a middle $\tau$ of a constant has one $\sigma$ before it
    inside the constant, and $W$'s head $\sigma^2$ must fit in the
    preceding $\sigma$-run. If that constant had more than $k$ $\tau$'s,
    the $(k{+}1)$-st would follow a gap of $1$ where $W$ has its seam
    $3$ or $5$: impossible; so the constant has exactly $k$ $\tau$'s,
    and by the distinctness of family indices it is the unique slot
    holding that string. The same argument applies constant-by-constant
    along $W$'s seams (a seam of $5$ matches a junction, so the
    constants on both sides occupy adjacent slots). Assembling:
    $M P_i M$ can only chain the $M$-slot, a junction, the $P_i$-slot,
    a junction, the $M$-slot, and the two $M$-slots of a configuration
    bracket exactly one code -- the wrapped one: (ii). $D_r D_{r'}$
    chains $D_r$'s slot to the next only through a gap of $5$, which
    after $D_r$'s slot arises exactly when the tally is empty: (iii).
    The output prefix is pinned at $D_0$'s slot and then slot-by-slot
    through the program zone (matching $\Pi(0)$ iff $i = 0$), then
    across the first tally junction ($5$ iff $x = 0$): (v). Finally a
    single frame constant $V(k)$ is itself such a pattern, pinned by
    Step 3 to a slot holding that string: the two $M$-slots are the two
    wrap positions, every other constant has one slot: (i).
\end{proof}

\begin{proof}[Proof of Theorem~\ref{thm:lim2cm} (correctness)]
    $\mathrm{INIT} = D_0\,\Pi(1)\,D_1\cdot[\,\tau/\sigma\,]X_1\cdot D_2 D_3$
    maps any $w$ over $\{\sigma, \tau\}$ to $C(1, |w|, 0)$: one constant
    pass turns the input into a $\tau$-tally of its length, and the
    frame constants sandwich it. By Lemma~\ref{lem:occurrence} the
    dispatch test of the current instruction is the only one that
    fires, the operation pass acts on exactly the intended tally, and
    unwrap-then-wrap moves the wrap to exactly the target code -- so an
    induction along the run of $M$ gives
    $\mathrm{step}(C(i, x, y)) = C(i', x', y')$ for the machine's step,
    and the halted configuration -- where every dispatch test fails and
    the identity branch runs -- is a fixed point. It is the unique
    fixed point among configurations: every non-halted configuration
    moves, since its instruction either edits a tally or jumps (the
    zero-branch self-jump that would not move is excluded by (N1)).
    The orbit of $\mathrm{lim}(\mathrm{step})$ on a configuration is
    therefore exactly the run of $M$, sweep for step: it converges iff
    $M$ halts, in as many sweeps as the machine takes steps. At halt,
    (N2) puts the result in counter 2 and (N3) empties counter 1, so
    Lemma~\ref{lem:occurrence}(v) applies to the configuration
    $C(0, 0, y)$: $\mathrm{OUT} =
    [\,\epsilon/D_0\Pi(0)D_1D_2\,][\,\epsilon/D_3\,]$ deletes the
    constant prefix (which occurs exactly once, at position $0$) and
    the closing delimiter, leaving $\tau^{y}$. Both deletion passes sit
    outside the iterand, so $\mathrm{step}$ itself is total.

    Verified end to end (the scripts and full domains are in the
    research log, research/scratch/lim/): the compiled doubling machine
    on $n \leq 12$ ($\tau^{2n}$ exact), the adder on all $25$ input
    pairs $n, m \leq 4$ plus a $7 \times 7$ border grid, a
    counter-draining machine on a $5 \times 5$ grid, and the
    normalization pipeline (N1)--(N3) demonstrated on a machine
    needing all three, with the raw defects shown (a non-halting fixed
    point, a counter-1 output, a dirty halt) and repaired; divergence
    verified on a self-looping and a two-cycling machine, the orbit
    matching the simulated run on every point; $\mathrm{step}$'s
    totality on $400$ random large strings and on every branch; the
    occurrence census of Lemma~\ref{lem:occurrence} at the edge
    lengths $x, y \in \{0, 1, 2\}$ for every instruction counter,
    systematically for $x, y \leq 40$, and on $400$ random
    configurations up to $1000$, at three constant assignments (a
    fourth, violating the distinctness hypothesis, run as a stress
    test); and the whole of it re-run under both swaps of the
    Section-2 toolkit roles.
\end{proof}

\begin{theorem}
    (\emph{The Compiled Program Computes the Machine})
    \label{thm:limpartial}
    Over any alphabet with $|\Sigma| \geq 2$, the flat calculus with one
    $\mathrm{lim}$ node computes exactly the partial computable
    functions $\Sigma^* \rightharpoonup \Sigma^*$: it computes all of
    them (Theorem~\ref{thm:lim2cm} and Minsky's two-counter machines
    \cite{minsky67}, arbitrary inputs brought to tallies by the
    two-letter prefix-code reduction of Theorem~\ref{thm:universal}),
    and it computes nothing else: a converging evaluation is a finite
    sequence of sweeps, each a primitive-recursive pass, halted by a
    decidable fixed-point test, so every denotation is partial
    computable.
\end{theorem}

\begin{proposition}
    (\emph{One Node Is Enough: the Driver})
    \label{prop:limcontain}
    Every $L{+}\mathrm{lim}$ expression is computed by a lazy-pass
    recursive $L$ expression of the same size up to a constant factor:
    the unary function $\mathrm{lim}(E)$ over a flat iterand $E$ is the
    guarded definition
    \begin{align}
        \mathrm{RUN}(X) \;\triangleq\; \mathtt{if}\big(\mathtt{eq}(E(X),\, X),\;\; X,\;\; \mathrm{RUN}(E(X))\big),
    \end{align}
    with the recursive call in the \emph{dead} branch's replacement --
    the pattern-gated position of Remark~\ref{rem:ifgate}, where the
    two-way gate of Theorem~\ref{thm:gate} makes $\mathtt{if}$ sound
    both ways.
\end{proposition}

\begin{proof}
    Induction on the expression, replacing each $\mathrm{lim}$ node by
    its driver bottom-up: the iterand $E$ is then $\mathrm{lim}$-free
    by the induction hypothesis, and the node
    $\mathrm{lim}(E)[E_0]$ is the driver applied to the start,
    $\mathrm{RUN}[E_0/X]$. The driver unrolls the orbit exactly: its
    guard tests $s_{k+1} = s_k$, and its else-branch forces $E(X)$ --
    one sweep -- before recurring, so the recursion halts on the least
    $K$ with $s_{K+1} = s_K$ and diverges with the orbit; and if
    $E(X)$ is undefined at any point, the guard is undefined and so is
    the driver, exactly as the orbit's undefinedness makes the node
    undefined (Definition~\ref{def:lim}) -- the two are undefined
    together, as in Lemma~\ref{lem:beta}'s caveat. What the eager
    calculus could not do -- evaluate a branch that may diverge -- is
    not asked of it: the dead branch's round is a discarded
    replacement (Remark~\ref{rem:ifgate}'s minimal repair), and the
    two-way gate of Theorem~\ref{thm:gate} is what makes that sound.
    Verified: the driver against the direct orbit for eight rules
    (convergent and divergent, census witnesses included) on $19$
    inputs each under both choices of the toolkit roles -- $135$
    convergent orbits value-exact, $29$ divergent -- plus the compiled
    iterands of Theorem~\ref{thm:lim2cm} (the drivers of the doubling
    and adder machines match $\mathrm{MAIN}$ exactly on $n \leq 5$) and
    a self-looping machine diverging on both sides.
\end{proof}

\begin{corollary}
    (\emph{One Node Is the Whole Gap})
    \label{cor:limclass}
    Over any alphabet with $|\Sigma| \geq 2$:
    \begin{align}
        \text{flat } L + \text{ one } \mathrm{lim} \text{ node}
        \;&=\; \text{lazy-pass recursive } L \notag \\
        \;&=\; \text{the partial computable functions.}
    \end{align}
\end{corollary}

\begin{proof}
    Left to right by Proposition~\ref{prop:limcontain} (one driver per
    node; the calculus underneath is untouched), right to left by
    Theorem~\ref{thm:limpartial} together with the soundness direction
    of Theorem~\ref{thm:universal}.
\end{proof}

This is the inertness story of Section~\ref{sec:recursion}, sharpened
to a single quantifier. The bounded calculus is polynomial
(Theorem~\ref{thm:fp}); recursion under eager evaluation changes
nothing (Theorem~\ref{thm:eager}), and neither do lazy arguments
(Theorem~\ref{thm:lazyargs}); the two-way gate plus recursion is
universal (Theorem~\ref{thm:universal}). The gate's condition --
evaluate the replacement only when the pattern fires -- is exactly the
fixed-point test of Definition~\ref{def:lim}, and
Proposition~\ref{prop:limcontain} says the translation is an
equality of classes, not an embedding of one into a richer one: the
power was never in the recursion depth, it is unbounded guarded
iteration, and one iteration node over a flat pipeline already has
all of it.

\begin{corollary}
    (\emph{Undecidability})
    \label{cor:limundec}
    \begin{enumerate}[(i)]
        \item Given a flat $E \in \mathrm{Exp}_1$ and $w$, whether
        $\mathrm{lim}(E)(w)$ is defined is $\Sigma^0_1$-complete.
        \item Given a flat $E \in \mathrm{Exp}_1$, whether
        $\mathrm{lim}(E)$ is total is $\Pi^0_2$-complete.
    \end{enumerate}
\end{corollary}

\begin{proof}
    Both problems are in their classes because ``the orbit converges
    within $k$ sweeps'' is decidable, each sweep a primitive-recursive
    pass. Hardness: $\mathrm{MAIN}(\tau^{n})$ is defined iff $M$ halts
    on $n$ (Theorem~\ref{thm:lim2cm}), a $\Sigma^0_1$-hard relation,
    and $\mathrm{MAIN}$ is total iff $M$ halts on \emph{every} input
    (every $n$ is $\mathrm{INIT}$'s image of some string -- every
    length is hit), a $\Pi^0_2$-hard one. The undecidability of
    totality is the sibling of Corollary~\ref{cor:barriers}(4) for the
    recursive calculus.
\end{proof}

\begin{proposition}
    (\emph{Growth})
    \label{prop:limgrowth}
    \begin{enumerate}[(i)]
        \item (\emph{exponential output, one node}) Writing $v$ for the
        Horner value of Theorem~\ref{thm:amplifier} ($a$ adds $1$, $b$
        doubles),
        \begin{align}
            \mathrm{lim}([\,baa/ab\,])(w) \;=\; b^{\,\#b(w)}\, a^{\,v(w)}
        \end{align}
        for every $w$, in at most $v(w) - \#a(w)$ sweeps; in particular
        \begin{align}
            \mathrm{lim}([\,baa/ab\,])(a\,b^{\,n-1}) \;=\;
            b^{\,n-1} a^{\,2^{\,n-1}} \text{ of length } 2^{\,n-1} + n - 1 .
        \end{align}
        One constant pass under one $\mathrm{lim}$ node has exponential
        output: the length bound of Lemma~\ref{lem:length} fails at
        depth one.
        \item (\emph{towers, nested nodes}) The block
        $\mathrm{lim}([\,baa/ab\,]) \circ \mathrm{lim}([\,ab/aa\,])$
        maps $b^{m} a^{K}$ to
        $b^{\,m + \lfloor K/2 \rfloor}\,
        a^{\,2^{\lfloor K/2 \rfloor + 1} - 2 + (K \bmod 2)}$ -- the
        closed form of Corollary~\ref{cor:towers}'s two restart nodes,
        now for sweep normalizers -- so $2t - 1$ constant-pattern
        $\mathrm{lim}$ nodes produce output of tower height $t$ in the
        input length, and $\mathrm{lim}(L) \sqsubseteq L$ fails: three
        nodes already outgrow Lemma~\ref{lem:length}'s polynomial
        bound.
    \end{enumerate}
\end{proposition}

\begin{proof}
    (i) Each replacement $ab \to baa$ preserves $\#b$ and $v$:
    reading $v$ left to right, $v(x\,ab) = (v(x)+1) \cdot 2 =
    v(x\,baa)$, and appending the same suffix multiplies both equally
    -- formally $v(uv) = v(u) \cdot 2^{\#b(v)} + v(v)$. So the sweep,
    which performs several such replacements at disjoint sites,
    preserves them too. The potential $\Phi(w) = v(w) - \#a(w)$ is
    nonnegative -- $v$ sums $2^{\#b(\text{after})}$ over the $a$'s of
    $w$ -- drops by exactly $1$ per replacement ($\#a$ grows, $v$
    does not), and is $0$ exactly when no $a$ has a $b$ after it, i.e.\
    on the $ab$-free strings, which by Lemma~\ref{lem:fpsweep} are the
    sweep's fixed points. The orbit therefore reaches an $ab$-free
    string after at most $\Phi(w)$ replacements -- hence at most that
    many sweeps -- and an $ab$-free string with the same $\#b$ and $v$
    is $b^{\#b} a^{v}$, both invariants forcing the exponents. The
    potential also re-proves the termination of the restart process of
    Theorem~\ref{thm:amplifier}, and the two normalizers of this rule
    agree on every input. (Verified: all binary inputs of length
    $\leq 12$, $200$ random of length $13$--$16$, the invariants per
    sweep on $50$ traces, and the tally witness for $n \leq 14$.)
    (ii) The half node maps $b^{m} a^{K}$ to
    $b^{m} (ab)^{\lfloor K/2 \rfloor} a^{K \bmod 2}$ in \emph{one}
    sweep -- the greedy pairing of the $a$-run -- and is then inert
    ($ab$ contains no $aa$); composing with (i) and counting the
    $b$'s after each $a$ of $(ab)^{j}$ gives the displayed form, the
    Horner computation of Corollary~\ref{cor:towers}. (Verified
    against the closed forms for $t \leq 3$ -- $5$ nodes, output
    length $65{,}556$ at $n = 4$ -- and the block on the grid
    $m \leq 3$, $K \leq 14$.)
\end{proof}

\begin{remark}
    (\emph{Divergent Orbits Grow Like Fibonacci})
    \label{rem:fiborbit}
    The census's divergent rule $[\,aabb/ba\,]$ diverges under the sweep
    too, but with a law: on $aabbaabb$ the orbit lengths are
    $2F_{k+1} + 2k + 6$ for the Fibonacci numbers $F$, i.e.\ $8 \to 10
    \to 14 \to 18 \to \cdots \to 204$ in ten sweeps $\to \cdots \to
    1{,}028{,}520$ in twenty-eight -- growth by the golden ratio per
    sweep, exponential in the iteration count; on $bbaa$ the law is
    $2F_{k+1} + 2k + 2$, and the mirror rule $[\,bbaa/ab\,]$ runs the
    reversed orbits (verified exactly through $k = 28$, the sweep
    strings cross-checked against the evaluator through $k = 16$). The
    mechanism is the census splice of
    Theorem~\ref{thm:termination}(v) -- a prefix ending $aab$ against
    the replacement $aabb$, re-creating a junction $ba$ -- compounding:
    a junction in the context $(b, a)$ doubles per sweep, a junction in
    $(a, a)$ or $(b, b)$ persists, one in $(a, b)$ dies, and the two
    persistent species follow the Fibonacci recurrence. So the growth
    of the iterated calculus lives in the \emph{iteration}, not the
    nesting: Corollary~\ref{cor:towers} needs $2t-1$ nested nodes for
    tower height $t$, but one node already produces exponential output
    (Proposition~\ref{prop:limgrowth}(i)) and Fibonacci divergent
    orbits -- while divergent restart rules were the \emph{only}
    obstacle left in the one-rule census, here divergence is the
    engine.
\end{remark}

The landscape closes one row and opens a question. The table of
Section~\ref{sec:variants} gains:

\begin{center}
\footnotesize
\begin{tabular}{l|llll}
    variant & $L \sqsubseteq V$ & $V \sqsubseteq L$ & toolkit & growth \\
    \hline
    $\mathrm{lim}$ & trivially & refuted & flat & unbounded
\end{tabular}
\end{center}

\noindent $L \sqsubseteq \mathrm{lim}(L)$ trivially, the grammar growing
by one node kind; $\mathrm{lim}(L) \sqsubseteq L$ refuted by
Proposition~\ref{prop:limgrowth} (one node, exponential output); the
toolkit carried verbatim -- the compilation of
Theorem~\ref{thm:lim2cm} needs no recursion and no nesting, an if-chain
of Section~\ref{sec:toolkit} passes; growth unbounded, the class being
the partial computable functions. Restart and $\mathrm{lim}$ are the
two iteration operators over the same grammar -- one rule iterated
against one expression iterated -- and both escape $L$ on growth
(Corollary~\ref{cor:towers}, Proposition~\ref{prop:limgrowth}), but
$\mathrm{lim}$'s class is exactly characterized
(Corollary~\ref{cor:limclass}) where the Markov row's entries carry
Conjecture~\ref{conj:ic}. What the characterization does \emph{not}
settle is the one-rule fragment: is convergence of
$\mathrm{lim}([\,A/B\,])$ decidable? Proposition~\ref{prop:limtotal}
settles everything outside the strategic case $|A| > |B|$ with
$B \not\subset A$, and Remark~\ref{rem:limstrategies} shows the sweep
and leftmost versions of that case are genuinely different problems --
the sweep sibling of Theorem~\ref{thm:termination}(v)'s strategic
termination problem, whose unrestricted form is the long-standing open
one-rule termination problem \cite{geser01}. And the hierarchy inside the
operator: one $\mathrm{lim}$ node over a \emph{single constant pass}
reaches exponential output and Fibonacci orbits
(Proposition~\ref{prop:limgrowth}, Remark~\ref{rem:fiborbit}) -- whether
it reaches universality is open; the amplifier of
Proposition~\ref{prop:limgrowth}(i) sorts and Horner-codes, and every
fixed point of a one-rule sweep is $B$-free
(Lemma~\ref{lem:fpsweep}), but no argument is known in either
direction.
